%%%%%%%%%%%%%%%%%%%%%%%%%%%%%%%%%%%%%%%%%%%%%%
% Krishna's macros for conf papers
%%%%%%%%%%%%%%%%%%%%%%%%%%%%%%%%%%%%%%%%%%%%%%

% Macros for reviewing paper, marking todos
\newcommand{\todo}[1]{{\color{red} TODO: \textbf{#1}}}
% \newcommand{\todo}[1]{}
% For comments
\newcommand{\krishna}[1]{{\color{red}{\textbf{Krishna}: #1}}}
\newcommand{\liam}[1]{{\color{blue}{\textbf{Liam}: #1}}}
\newcommand{\qiao}[1]{ {\color{cyan}{\textbf{Qiao}: #1}} }
\newcommand{\shuang}[1]{{\color{green}{\textbf{shuang}: #1}}}

% For highlighting paper edits
\newcommand{\edit}[1]{{\color{red}{#1}}}

% Color defined by Florian Golemo: https://fgolemo.github.io/
% Good choice for references
\definecolor{flodarkpurple}{rgb}{0.288,0.1196,0.7}

\definecolor{ganeshamber}{rgb}{1.0, 0.75, 0.0}
\newcommand{\ganesh}[1]{{\color{ganeshamber} \textbf{Ganesh}: \textbf{#1}}}



% Shorthand for \mathbf
\newcommand{\mb}[1]{\mathbf{#1}}
% Shorthand for \mathcal
\newcommand{\mc}[1]{\mathcal{#1}}
% Short hand for delta * t (Defined by: Miles Macklin)
\newcommand{\dt}{\Delta t}

% % Common abbreviations
\newcommand{\etal}{\emph{et al.}}
\newcommand{\cf}{\emph{cf.}}


% Cross mark (assumes package `pifont` is used)
\newcommand{\xmark}{\ding{55}}%

% % Table float box with bottom caption, box width adjusted to content
% % (assumes package `floatrow` is used)
% \newfloatcommand{capbtabbox}{table}[][\FBwidth]

% % Text colors
% \newcommand{\red}[1]{{\textcolor{red}{#1}}}
% \newcommand{\blue}[1]{{\textcolor{blue}{#1}}}
% \newcommand{\green}[1]{{\textcolor{green}{#1}}}
% \newcommand{\gray}[1]{{\textcolor{gray}{#1}}}
% \newcommand{\pink}[1]{{\textcolor{pink}{#1}}}

% % % Credits: Derek Nowrouzezahrai
% % \newenvironment{myitemize}{
% % \begin{itemize}[leftmargin=*,noitemsep,topsep=0pt,parsep=0pt,partopsep=0pt]
% %   \setlength{\itemsep}{1pt}
% %   \setlength{\parskip}{1pt}
% %   \setlength{\parsep}{1pt}}{\end{itemize}
% % }

% % % Credits: Derek Nowrouzezahrai
% % \newenvironment{myenumerate}{
% % \begin{enumerate}[leftmargin=*,noitemsep,topsep=2pt,parsep=0pt,partopsep=0pt]
% % %  \setlength{\parindent}{0pt}
% %   \setlength{\itemsep}{1pt}
% % %  \setlength{\itemindent}{6pt}
% %   \setlength{\parskip}{1pt}
% % %  \setlength{\parsep}{0pt}
% % }{\end{enumerate}
% % }
